\scanpage{118}
\deleted{\fbox{1''}} Quadriques affines \inserted{bi}régulières (i.e. \inserted{\unsure{project\textsuperscript{t}|projectivement}} lisses et lisses à l'infini \unsure{voir --|non --})

1. Soient, sur schéma $S$, $P$ un fibré projectif de rang $2\nu-1$, $Q$ une quadrique lisse dans $P$, $H$ un hyperplan de $P$ non tangent à $Q$, $\unsure{P|Q}-H=E$ le fibré affine correspondant. Donc $E$ (ou encore $(P,H)$) est déterminé par une extension de $\mathcal O_S$-fibrés vect.
\diagram{\[
\begin{tikzcd}[ampersand replacement=\&,column sep=large]
0 \arrow[r] \& V \arrow[r] \& T \arrow[r,"\pi"] \& \mathcal O_S \arrow[r,phantom,"{\unsure{-|\longrightarrow}}" description] \& 0
\end{tikzcd}\tag*{(1)}
\]}
où $\underline V$ est le $\mathcal O_S$-Module loc. libre (de rang $2\nu-1$) des translations de $T$, $E=\pi^{-1}(1)$, \inserted{$T$ étant de rg $2\nu$.} Quant à $Q$, elle est déterminée par la donnée d'une section de $\check{\mathbb P}(\operatorname{Sym}^2(\check T))=\mathbb P(\deleted{\unsure{\operatorname{Sym}}}\Gamma^2(T))$, i.e. une forme quadratique \inserted{$q$} sur $T$ définie loc. mod. multiplication par une section inversible.

La condition $Q$ lisse signifie $q$ lisse i.e. la forme bil. associée $\varphi$ non dégénérée, ou encore que
\[
\delta(q)\in\Gamma(S,(\det T)^{\otimes-2})\tag*{(2)}
\]
est une \deleted{section}\inserted{base} de $(\underline{\det}\,T)^{\otimes-2}$, i.e. est partout non \unsure{nulle}. La condition que $H$ est non tangent à $Q$ signifie aussi que $q\mid V$ est une forme lisse. \deleted{ou encore que la droite ortho}\inserted{On va la} reformuler, en notant $D$ le \unsure{$k$-fibré|fibré|ss-fibré|sous-fibré} inversible
% archive sheet number: 117
